% Notes-only proposal. Root owns integration, numbering and register changes.
% The three corrections are auxiliary source audits, not movement of either
% outstanding subject-21 original-source gate.
\paragraph{Requested determinant background: conventions and scope.}
The original of cote~68, scan~056, concerns finite free modules,
free rank-one graded modules with Koszul commutativity, and invariance
under changing a free complement.  It is a finite-free prototype of the
determinant formalism.  The bounded-complex and perfect-complex data in
Knudsen--Mumford's Definitions~1 and~4 have additional requirements
\cite[pp.~20--24, 40--42]{ROOT_KnudsenMumford1976}.
The following proofs also record three qualifications found by reading
all of the original article, pp.~19--55.

\begin{lemma}[Proved: graded signs, exterior blocks and the printed sum-exponent defect]\label{C:graded-determinant-sign}
Let $S$ be any scheme.  Let $\operatorname{Pic}^{\mathbb Z}(S)$ be the
groupoid whose objects are pairs $(L,a)$, with $L$ an invertible
$\mathcal O_S$-module and $a:S\to\mathbb Z$ locally constant, and whose
arrows are degree-preserving $\mathcal O_S$-linear isomorphisms.  Set
\[
 (L,a)\otimes(M,b)=(L\otimes M,a+b),\qquad
 \symbf{1}=(\mathcal O_S,0).
\]
With the ordinary tensor associator and unit maps, the formula
\[
 c_{(L,a),(M,b)}(l\otimes m)=(-1)^{ab}m\otimes l
\]
defines a symmetric monoidal Picard groupoid.  The functor
\[
 E\longmapsto\operatorname{Det}(E)
       =(\textstyle\bigwedge^{\operatorname{rk}E}E,\operatorname{rk}E)
\]
on finite locally free modules and their isomorphisms, equipped with
the exterior-product maps for direct sums, respects this symmetry.
It respects all finite block permutations and, on finite free modules,
the change-of-complement condition of scan~056.

The replacement $(-1)^{a+b}$ printed in the commutativity formula of
Knudsen--Mumford, p.~20, fails these unit and exterior-product
requirements already over $\operatorname{Spec}\mathbb Q$.
\end{lemma}
\begin{proof}
The function $(-1)^{ab}$ is a locally constant unit of
$\mathcal O_S$, so the displayed formula is a globally defined line
isomorphism.  It is natural because line morphisms are linear.
The composite of the two swaps has scalar
$(-1)^{ab+ba}=1$.  Under the ordinary tensor associator the two
hexagon identities reduce, respectively, to
\[
 a(b+c)=ab+ac,\qquad (a+b)c=ac+bc.
\]
The unit compatibility follows from $a\cdot0=0$.
The associator has its ordinary tensor pentagon, and
$(L^\vee,-a)$ is a tensor inverse of $(L,a)$, by evaluation.
This verifies the stated Picard structure for arbitrary integral
degrees, including negative ones.

All assertions about finite locally free modules can be checked where
their ranks are constant and bases exist.  For bases
$e_1,\ldots,e_r$ of $E$ and $f_1,\ldots,f_s$ of $F$, let
$i_{E,F}$ send the tensor of the two volume forms to their concatenated
wedge in $\bigwedge^{r+s}(E\oplus F)$.  Let
$\tau:E\oplus F\to F\oplus E$ send $(e,f)$ to $(f,e)$.
Moving the $r$ vectors of $E$ past the $s$ vectors of $F$ makes exactly
$rs$ elementary exterior swaps.  Consequently
\[
 \operatorname{Det}(\tau)\,i_{E,F}
       =i_{F,E}\,c_{\operatorname{Det}(E),\operatorname{Det}(F)}.
\]
Associativity of the maps $i$ is associativity of concatenation;
their unit normalization uses the empty wedge.  For a permutation
which sends block $i$ to position $\sigma(i)$, every inverted pair
$i<j$ with $\sigma(i)>\sigma(j)$ contributes $r_i r_j$ crossings.
Both the exterior and graded-tensor routes therefore have the sign
\[
 (-1)^{\sum_{i<j,\ \sigma(i)>\sigma(j)}r_i r_j}.
\]
This proves the assertion for every finite block permutation.

For the manuscript's complement condition, write $V=W\oplus U$
and $V=W\oplus U'$ with all indicated modules finite free.
Let $\lambda:U\to U'$ be the isomorphism inducing the identity
on $V/W$.  For a basis $w_1,\ldots,w_r$ of $W$ and a basis
$u_1,\ldots,u_s$ of $U$, put $u'_j=\lambda(u_j)$.
Then $u'_j-u_j\in W$.  In
\[
 w_1\wedge\cdots\wedge w_r\wedge u'_1\wedge\cdots\wedge u'_s
\]
every expansion term other than the one with all the $u_j$ vanishes,
since it contains $r+1$ vectors from $W$.  The two determinant
identifications are equal.  Equivalently, an automorphism which is
the identity on $W$ and $V/W$ has block form
$\left(\begin{smallmatrix}I&T\\0&I\end{smallmatrix}\right)$,
and the same wedge calculation gives determinant one.  The rank-one
normalization is $\bigwedge^1L=L$ in degree one.

Finally, over $\mathbb Q$, degrees $(1,0)$ give scalar $-1$ under
the printed sum rule, although swapping with the specified unit has
scalar $+1$.  For two rank-one blocks, degrees $(1,1)$ give $+1$
under that rule, whereas $f\wedge e=-e\wedge f$.  Thus both are
actual counterexamples in characteristic zero.  The sum rule is
nevertheless involutive, so the double-swap identity alone would not
detect its defect.
\end{proof}

Grothendieck's letter to Knudsen dated 19 May 1973, as reproduced in
the opened Appendix~B of Knudsen's later paper, explicitly uses the
product of the two degrees and requires associativity and unit
compatibility \cite[Appendix~B, author-version pp.~40--42, especially
p.~41]{A567_GrothendieckKnudsen1973Letter}.
This supplies original-context corroboration; it is not a claim to
have located a formal erratum to the 1976 printing.
The exact finite certificate \texttt{C\_graded\_determinant\_sign\_certificate.py}
checks four parity pairs, eight parity triples and the 25 block-rank
pairs with ranks between zero and four.  Its 90 assertions passed.
The general assertions are established by the proof above, not by
extrapolation from those finite cases.

\begin{proposition}[Refuted: unrestricted ordinary base change in the printed pushforward paragraph]\label{C:determinant-ordinary-basechange}
The assumptions ``$p:X\to Y$ proper of finite Tor dimension,
$Y$ noetherian, $F$ perfect on $X$, and $g:Y'\to Y$ a morphism
of noetherian schemes'' do not imply the ordinary-fibre-product
comparison displayed on p.~46 of
\cite{ROOT_KnudsenMumford1976}.  With
$X'=X\times_Y Y'$ and projections $g':X'\to X$, $p':X'\to Y'$,
the two graded lines
\[
 \operatorname{Det}(Lg^*Rp_*F)
 \quad\hbox{and}\quad
 \operatorname{Det}(Rp'_*Lg'^*F)
\]
need not even have the same degree.
\end{proposition}
\begin{proof}
Put $R=\mathbb Q[t]$, $k=R/(t)$,
$Y=\operatorname{Spec}R$ and
$X=Y'=\operatorname{Spec}k$.  Let both $p$ and $g$ be the closed
immersion induced by $R\to k$, and set $F=\mathcal O_X$ in degree
zero.  The morphism $p$ is proper, and all schemes are noetherian.
The exact free resolution
\[
 0\longrightarrow R\xrightarrow{\ t\ }R\longrightarrow k
   \longrightarrow0
\]
shows that $p$ has finite Tor dimension one.  The complex $F$ is
perfect on $X$.  The ordinary tensor product $k\otimes_R k=k$
identifies $X'$ with $\operatorname{Spec}k$; both $g'$ and $p'$
are identity maps.

Since a closed immersion is affine, $Rp_*F$ is the $R$-module $k$
in degree zero.  Pulling its displayed free resolution back gives
\[
 Lg^*Rp_*F=k\otimes_R^{\mathbf L}k
       = [\,k\xrightarrow{\ 0\ }k\,],
       \qquad\text{in degrees }-1,0.
\]
Its graded determinant has degree $-1+1=0$.  In contrast,
$Rp'_*Lg'^*F=k$ in degree zero has determinant degree one.
Degree-preserving line isomorphisms cannot identify these objects.
The derived comparison is already not a quasi-isomorphism, since
its source has $H^{-1}=k$ and its target has $H^{-1}=0$.
The obstruction is the nonzero
$\operatorname{Tor}_1^R(k,k)=k$; the square is not Tor independent.
The first isomorphism printed on p.~46, between the pullback of a
determinant line and the determinant of the derived pullback, is not
challenged by this example.  It is the second, pushforward base-change
comparison that fails.
\end{proof}

\begin{proposition}[Proved: the sharp affine criterion and Tor-independent repair]\label{C:determinant-basechange-repair}
Let $R\to S$ and $R\to A$ be homomorphisms of commutative rings,
and put $B=A\otimes_R S$.  For a perfect $S$-complex $F$, consider
the classical comparison in $D(A)$
\[
 \beta_F:A\otimes_R^{\mathbf L}F
       \longrightarrow B\otimes_S^{\mathbf L}F,
\]
where the target is restricted along $A\to B$.
Then $\beta_F$ is a quasi-isomorphism for every perfect $S$-complex
$F$ if and only if $\operatorname{Tor}_i^R(A,S)=0$ for all $i>0$.

More generally, let $p:X\to Y$ be a quasi-compact, quasi-separated
morphism of schemes, let $g:Y'\to Y$ be any morphism, and form
$X'=X\times_Y Y'$ with projections $g'$ and $p'$.  Suppose that
$X$ and $Y'$ are Tor independent over $Y$: on affine charts over
an affine open of $Y$, all higher Tor modules of their two coordinate
algebras vanish.  For every
$E\in D_{\mathrm{QCoh}}(\mathcal O_X)$ the canonical map
\[
 Lg^*Rp_*E\longrightarrow Rp'_*Lg'^*E
\]
is an isomorphism.  If $Rp_*E$ is perfect, this gives an isomorphism
of their graded determinants as well.  In particular, flatness of
either $p$ or $g$ supplies the Tor-independence hypothesis.
When $p$ is proper of finite Tor dimension, $Y$ is noetherian and
$E$ is perfect, $Rp_*E$ is perfect, so this is a corrected version
of the printed determinant comparison.
\end{proposition}
\begin{proof}
For necessity in the affine assertion take $F=S[0]$.
The target is $B[0]$, whereas the source has
$H^{-i}=\operatorname{Tor}_i^R(A,S)$ for $i\geq0$.
For sufficiency choose a free $R$-module resolution $P^\bullet\to A$
in nonpositive degrees and represent $F$ by a bounded complex
$G^\bullet$ of finite projective $S$-modules.  After forgetting the
$A$-action, the canonical comparison is represented by
\[
 (P^\bullet\otimes_R S)\otimes_S G^\bullet
       \longrightarrow B\otimes_S G^\bullet.
\]
The assumed Tor vanishing makes $P^\bullet\otimes_R S\to B$
a quasi-isomorphism of $S$-complexes.  Tensoring with a single
projective $S$-module preserves an acyclic complex.  The finite
filtration of $G^\bullet$ by its degrees therefore shows, successively,
that tensoring with $G^\bullet$ preserves this quasi-isomorphism.
Forgetting from $D(A)$ to $D(R)$ detects quasi-isomorphisms, since
it does not change the underlying cohomology groups.  This proves
sufficiency without incorrectly assigning an $A$-module structure
to an arbitrary chosen $R$-free resolution.

For the scheme assertion use the canonical base-change comparison
and cover $Y'$ by affine opens lying over affine opens of $Y$.
Restriction to these opens commutes with the sheaf-theoretic
derived images in question.  It suffices to take $g$ between affine schemes.
Then $g$ and $g'$ are affine, and $Rg_*$ is conservative on complexes
with quasi-coherent cohomology.  Indeed, affine pushforward is exact
and faithful on quasi-coherent modules, so it detects each cohomology
module.  Write $K=g_*\mathcal O_{Y'}$.  Pushing the comparison forward
along $g$ identifies it with
\[
 Rp_*E\otimes_{\mathcal O_Y}^{\mathbf L}K
       \longrightarrow
 Rp_*\bigl(E\otimes_{\mathcal O_X}^{\mathbf L}
                          g'_*\mathcal O_{X'}\bigr).
\]
The ordinary affine fibre-product calculation gives
$g'_*\mathcal O_{X'}=p^*K$.  Tor independence gives
$Lp^*K\simeq p^*K$, as follows on an affine chart from the
vanishing of the higher homology of the derived tensor product.
The displayed map is thus the derived projection-formula
isomorphism for the quasi-compact, quasi-separated map $p$.
This is the derivation in the opened Stacks proofs
\cite[Tags~08IB and 08EU]{BStacks}.
The latter projection formula can be checked after taking $Y$
affine: its valid objects $K$ are closed under shifts, triangles
and arbitrary direct sums, since derived tensor and $Rp_*$ on
quasi-coherent complexes preserve those sums.  Free resolutions
and their truncation homotopy colimits generate $D(R)$ from $R$,
where the formula is the identity.  This proves it for the required
$K$ and completes the comparison.

Derived pullback preserves a bounded locally free local model, so
if $Rp_*E$ is perfect then the left-hand object is perfect; the
isomorphism makes the right-hand object perfect too.  Their
determinants may therefore be taken.  If one side of the square
is flat, its local algebra is flat over the base, so its higher
Tor modules with the other algebra vanish.

For the final proper case, work locally on the noetherian target.
A perfect complex on the proper source has bounded coherent
cohomology and finite Tor amplitude relative to the target:
use its bounded locally free models and the finite Tor dimension
of $p$, with a finite cover over the chosen affine base.
Proper coherent pushforward gives bounded coherent cohomology
for $Rp_*E$, hence pseudo-coherence.  For every quasi-coherent
module $M$ on the base, the projection formula identifies
$Rp_*E\otimes^{\mathbf L}M$ with
$Rp_*(E\otimes_{p^{-1}\mathcal O_Y}^{\mathbf L}p^{-1}M)$.
Relative Tor amplitude bounds the latter uniformly from below;
the bounded-above complex $Rp_*E$ supplies the upper bound.
Thus $Rp_*E$ has finite Tor amplitude.  Pseudo-coherence and finite
Tor amplitude imply perfection.  These are precisely the steps
in the opened perfect-pushforward proof
\cite[Tags~0B6G and 08EV]{BStacks}.
\end{proof}

\begin{lemma}[Proved: the exact scalar trace-splitting condition]\label{C:determinant-trace-splitting}
Let $R$ be a commutative ring, let $m\geq1$, and set
$\mathfrak s_m(R)=\ker(\operatorname{tr}:M_m(R)\to R)$.
The natural map
\[
 \phi:R\oplus\mathfrak s_m(R)\longrightarrow M_m(R),
       \qquad (a,T)\longmapsto aI_m+T,
\]
fits into the exact sequence
\[
 0\longrightarrow R[m]\longrightarrow
 R\oplus\mathfrak s_m(R)\xrightarrow{\phi} M_m(R)
 \xrightarrow{\operatorname{tr}\bmod m}R/mR\longrightarrow0,
\]
where $R[m]=\{a\in R:ma=0\}$ and the first map sends
$a$ to $(a,-aI_m)$.
It is an isomorphism if and only if $m$ is a unit in $R$.
In particular, the scalar/trace-zero direct-sum description on p.~52
of \cite{ROOT_KnudsenMumford1976} is false over the unrestricted
noetherian bases stated there.  For $R=\mathbb F_2$, $m=2$ there
is not even a lift of $1$ under trace which is invariant under
every change of basis.
\end{lemma}
\begin{proof}
The equation $aI_m+T=0$ forces $T=-aI_m$, and this has trace zero
exactly when $ma=0$.  This proves the kernel assertion and the
injectivity of the first map.  The final trace map is onto because
$\operatorname{tr}(E_{11})=1$.  A matrix is in its kernel exactly
when its trace is $ma$ for some $a\in R$; subtracting $aI_m$ then
makes it trace zero.  This proves exactness at $M_m(R)$.

If $m$ is a unit, the inverse of $\phi$ is
\[
 A\longmapsto
 \bigl(m^{-1}\operatorname{tr}(A),\,
       A-m^{-1}\operatorname{tr}(A)I_m\bigr).
\]
Conversely, if $\phi$ is surjective, express $E_{11}=aI_m+T$
with $\operatorname{tr}(T)=0$.  Taking traces gives $1=ma$,
so $m$ is a unit.  This proves the equivalence.

For the explicit noetherian example $R=\mathbb F_2$, $m=2$,
the nonzero identity matrix has trace zero.  Thus the scalar line
intersects the trace-zero subspace nontrivially; equivalently
$\phi(1,I_2)=0$ with $(1,I_2)\ne0$.

For the stronger invariance assertion put
\[
 J=\begin{pmatrix}0&1\\1&0\end{pmatrix},\qquad
 U=\begin{pmatrix}1&1\\0&1\end{pmatrix}.
\]
Both matrices are invertible over $\mathbb F_2$.
Writing $M=\left(\begin{smallmatrix}a&b\\c&d\end{smallmatrix}\right)$,
the equation $MJ=JM$ gives $c=b$ and $d=a$.
Substitution in $MU=UM$ then gives $b=0$.  Thus a matrix fixed by
conjugation by both $J$ and $U$ is scalar, and its trace is
$2a=0$.  Any change-of-basis-invariant lift of $1$ would commute
with both matrices and have trace one, a contradiction.
\end{proof}

\begin{lemma}[Proved: incidence and its Koszul resolution without trace splitting]\label{C:determinant-incidence}
Let $E$ be a finite locally free module of positive constant rank
$m$ on any scheme $S$.  Use the convention that $\mathbf P(E)$
parametrizes rank-one locally free quotients of $E$.
On $W=\mathbf P(E)\times_S\mathbf P(E^\vee)$ there is a canonical
section
\[
 \delta\in\Gamma(W,
  \operatorname{pr}_1^*\mathcal O_{\mathbf P(E)}(1)
  \otimes\operatorname{pr}_2^*\mathcal O_{\mathbf P(E^\vee)}(1))
\]
whose zero scheme is the incidence locus.  It commutes with arbitrary
base change, and
\[
 0\longrightarrow\mathcal O_W(-1,-1)
       \xrightarrow{\ \delta\ }\mathcal O_W
       \longrightarrow\mathcal O_{V(\delta)}\longrightarrow0
\]
is exact.  On a product with finitely many independent dual factors,
the tensor product of these two-term complexes resolves their
scheme-theoretic incidence intersection.  These assertions need
no invertibility condition on $m$.
\end{lemma}
\begin{proof}
For any finite locally free $\mathcal O_S$-module $E$, the canonical
isomorphism $E\otimes E^\vee\simeq\mathcal{E}nd(E)$ sends
$\sum_i e_i\otimes e_i^\vee$ to $\operatorname{id}_E$.
A basis change and its inverse dual cancel in this tensor, so it
is independent of the local basis.  Apply the two universal quotient
maps to this identity tensor.  This defines $\delta$, with local
equation $\sum_iX_iY_i$, and shows directly that it commutes with
base change.  On a test scheme, let the two quotients be
$E\twoheadrightarrow L$ and $E^\vee\twoheadrightarrow M$.
The section corresponds to the composite
$M^\vee\hookrightarrow E\twoheadrightarrow L$.
Its vanishing is precisely incidence, proving the description of
the zero scheme.

On $P=\mathbf P(E)$ the kernel $K$ of the universal quotient
$\pi^*E\twoheadrightarrow\mathcal O_P(1)$ is locally free of
rank $m-1$.  The quotient splits locally.  On such an open, choose
a basis with the quotient equal to the last-coordinate projection.
The incidence locus in the dual projective bundle is then
$\mathbf P(K^\vee)$, cut out by the last homogeneous coordinate.
On an affine projective chart the equation is either a polynomial
coordinate, which is a nonzerodivisor over any coefficient ring,
or a unit, in which case its zero scheme is empty.  The displayed
two-term complex is therefore exact on every chart and hence globally.

For several dual factors, take product affine charts after the same
local splitting over $P$.  If one incidence equation is a unit,
its two-term complex is contractible and the intersection is empty.
Otherwise the equations are distinct polynomial coordinates, one
from each factor.  Successive quotients remain polynomial rings
over the same coefficient ring, so each next coordinate is a
nonzerodivisor.  They form a regular sequence.  Inductively tensoring
the two-term free resolutions consequently resolves their common
quotient.  Exactness is local and gives the asserted global Koszul
resolution.  For $m=1$ each section is a unit and the same empty-locus
argument applies.
\end{proof}

Two further exact certificates supplement these proofs:
\begin{itemize}
\item \path{code/C_determinant_scope_certificate.py} records
35 passing assertions for the explicit derived-tensor dimensions and
the 16-matrix trace example.
\item \path{code/C_trace_invariance_certificate.py} records
22 passing assertions, including all 16 simultaneous-centralizer cases.
\end{itemize}
The certificates do not replace the general proofs above.

These corrections do not assert a failure of the determinant
construction.  In using its extended axioms one must retain the
original's strict diagrams of actual complexes: the footnote to
Definition~4 explicitly rules out replacing that condition by
arbitrary commutativity in the derived category.  The later extensions
to ordinary triangles have the additional perfect-cohomology or
reduced-base hypotheses stated there
\cite[pp.~37--46, especially p.~41 footnote and pp.~43--44]{ROOT_KnudsenMumford1976}.
